\begin{abstract}
Geometric spatial reasoning underlies applications ranging from robotics to automated design, yet the ability of language models to turn a verbal description into a physically connected 3D assembly remains largely unexplored. LEGO\textsuperscript{\tiny\textregistered}-brick assembly offers a demanding constructive test of this ability: a short, underspecified prompt leaves dimensions, proportions, and part layout implicit, while a valid assembly must both capture the prompt's semantic intent and satisfy the physical constraints imposed by interlocking studs. We formalize this setting as the \texttt{Text-to-Brick} task and present BrickForma, a framework that crystallizes the
geometric knowledge of large language models (LLMs) into small language models for this task, without the need of human-annotated
or human-verified 3D shape assets for training. BrickForma uses executable programs written by LLM agents to synthesize training data. These programs procedurally generate 79,301 3D voxel models across 21 categories on a $48\times48\times48$ grid. Each model is paired with five captions of increasing detail, yielding approximately 400,000 geometry-grounded training examples. We convert the voxel models into physically valid brick targets through connectivity analysis and refinement, and then fine-tune Qwen3-1.7B to predict brick dimensions and positions from text. Prior systems depend on curated assets: BrickGPT builds on human-annotated ShapeNetCore meshes, and BrickNet relies on human-designed LDraw brick assemblies. In contrast, BrickForma generates both the voxel geometries and their brick targets programmatically, removing the need for curated meshes and hand-built assemblies.
On 525 shared prompts spanning five description levels, BrickForma's outputs %before refinement 
achieve the highest VQAScore, PE, and SigLIP2 scores among the three systems at every level. Its invalid-placement ratio is 0.79\%, close to BrickGPT's 0.00\% and far below BrickNet's 18.42\%.
\end{abstract}